\documentclass[12pt]{article}
\usepackage[utf8]{inputenc}
\usepackage{amsmath}
\usepackage{amsthm}
\usepackage{amssymb}
\usepackage{graphicx}
\usepackage{natbib}
\usepackage{adjustbox}
\usepackage{fullpage}
\setlength{\parskip}{1em}
\title{Enumerative geometry of lines}
\author{Wei Jiang}
\date{August 2019}
\newcommand{\BigWedge}{\mathord{\adjustbox{valign=B,totalheight=.6\baselineskip}{$\bigwedge$}}}

\begin{document}



\theoremstyle{definition}
\newtheorem{df}{Definition}[section]
\newtheorem{eg}[df]{Example}


\theoremstyle{plain}
\newtheorem{thm}[df]{Thmorem}
\newtheorem{lm}[df]{Lemma}
\begin{titlepage}

\newcommand{\HRule}{\rule{\linewidth}{0.5mm}} % Defines a new command for the horizontal lines, change thickness here

%----------------------------------------------------------------------------------------
%	LOGO SECTION
%----------------------------------------------------------------------------------------

\includegraphics[width=8cm]{title/logo.eps}\\[1cm] % Include a department/university logo - this will require the graphicx package
 
%----------------------------------------------------------------------------------------

\center % Center everything on the page

%----------------------------------------------------------------------------------------
%	HEADING SECTIONS
%----------------------------------------------------------------------------------------

\textsc{\LARGE UROP Project}\\[1.5cm] % Name of your university/college
\textsc{\Large Imperial College London}\\[0.5cm] % Major heading such as course name
\textsc{\large Department of Mathematics}\\[0.5cm] % Minor heading such as course title

%----------------------------------------------------------------------------------------
%	TITLE SECTION
%----------------------------------------------------------------------------------------
\makeatletter
\HRule \\[0.4cm]
{ \huge \bfseries \@title}\\[0.4cm] % Title of your document
\HRule \\[1.5cm]
 
%----------------------------------------------------------------------------------------
%	AUTHOR SECTION
%----------------------------------------------------------------------------------------

\begin{minipage}{0.4\textwidth}
\begin{flushleft} \large
\emph{Author:}\\
\@author % Your name
\end{flushleft}
\end{minipage}
~
\begin{minipage}{0.4\textwidth}
\begin{flushright} \large
\emph{Supervisor:} \\
Dr. Arbesfeld Noah \\[1.2em] % Supervisor's Name

\end{flushright}
\end{minipage}\\[2cm]
\makeatother

% If you don't want a supervisor, uncomment the two lines below and remove the section above
%\Large \emph{Author:}\\
%John \textsc{Smith}\\[3cm] % Your name

%----------------------------------------------------------------------------------------
%	DATE SECTION
%----------------------------------------------------------------------------------------

{\large \today}\\[2cm] % Date, change the \today to a set date if you want to be precise

\vfill % Fill the rest of the page with whitespace

\end{titlepage}
\tableofcontents
\newpage

\section*{Introduction}
\newpage

\section{Grassmannian as projective variety}

How can we parametrized all lines in a projective space $\mathbb{P}^{n}$ ? We know a line $l$ in $\mathbb{P}^{n}$ corresponds to a 2-dimensional linear subspace of $\mathbb{A}^{n}$. So the the question transfers to parameterising linear subspaces of a vector space.



\begin{df}
    Let $V$ be a vector space of dimension $n$, define $G(k,n)$ to be the set of all linear $k$-subspaces of $V$.
\end{df}
 
\begin{eg}
    A first example of such a space is projective spaces $\mathbb{P}^{n}$ or they can be denoted as $G(1,n + 1)$. Lines in $\mathbb{P}^{n}$ form $G(2,n+1)$. 
\end{eg}

To avoid confusion, if $\Lambda$ is a k-dimensional subspace of a n-dimension vector space $V$, then we use $[\Lambda] \in G(k,n)$ to denote the point corresponds to the subspace.
\\

It can be shown that Grassmannian is a projective variety, so we can use all the tools in intersection theory to deal with it. The structure can be seen via a method called Plucker embedding.\\

Let $\Lambda \subset V$ be a k-dimensional subspace, then $\BigWedge^{k}\Lambda$ is a 1-dimensional subspace of $\BigWedge^{n}V$. More precisely, if $v_{1},...,v_{k}$ is a basis of $\Lambda$, then $\BigWedge^{k}\Lambda$ is the line spanned by $ v_{1}\wedge...\wedge v_{k}$. This corresponds to a point of $\mathbb{P}\BigWedge^{k}V$. \\

(By $\mathbb{P}V$ we mean the quotient of $V$ by the relation $v~w$ iff $v = \lambda w$. This is similar to the construction of projective space)\\

This gives a map of sets $G(k,n) \longrightarrow \mathbb{P}\BigWedge^{k}V \cong \mathbb{P}^{ {n \choose k} -1}$. This map is called Plucker embedding.

\begin{lm}
    This map  $G(k,n) \longrightarrow \mathbb{P}\BigWedge^{k}V \cong \mathbb{P}^{ {n \choose k} -1}$ is injective.
\end{lm}

\begin{proof}
    let $v_{1},...,v_{k}$ be a basis of $\Lambda \subset V$, we can extend it to a basis of $V$ by adding linear independent vectors $ u_{k+1},...u_{n}$. Let $a =  v_{1}\wedge...\wedge v_{k}$, then $\forall v \in V$, $v \wedge a = 0$ if and only if $ (b_{1}v_{1}+...b_{n}u_{n})\wedge a = 0$ iff  $b_{k+1}u_{k+1}\wedge a + ... +  b_{n}u_{n}\wedge a =0 $ iff $ v \in \Lambda$ so $ a$ determines $\Lambda$.
\end{proof}

So $G(k,n) $ is isomorphic to its image in $\mathbb{P}\BigWedge^{k}V$. Lets call its image $G$. It left to show that its image is the common zero locus of some homogeneous polynomials. We can show this by express our k-subspace as a matrix.\\


Let ${e_1,e_2,...,e_n}$ is a basis for $V$,we can identify $V$ as $k^{n}$ then any $k$-vector space is the span of $k$ linear independent vectors in this basis. We can write them as a matrix:
$$
\begin{pmatrix}
    a_{1,1}& a_{1,2} & \dots & a_{1,n}\\
    a_{2,1}& a_{2,2} & \dots & a_{2,n}\\
    \vdots & \vdots & \ddots & \vdots \\
    a_{k,1}& a_{k,2} & \dots & a_{k,n}    
\end{pmatrix}
$$



However, just like coordinates of points in a projective space, the matrix $A$ is not unique. Since we can multiply on the left any invertible $k \times k$ matrix $\Omega$ without changing the row spaces, because the rows in $\Omega A$ are linear combinations of the rows in $A$.  
In this setting, $\BigWedge^{k}V$ is given by the set


\[ \{e_{i1}\wedge...\wedge e_{ik}\}_{1\leq i1<...<ik\leq n}\]


After the Plucker embedding, this matrix get sent to the wedge product of row vectors:


\[ v_{1}\wedge...\wedge v_{k} = \sum_{1\leq i1<...<ik\leq n} D_{i1,...,ik} e_{i1}\wedge...\wedge e_{ik}   \]

Where $D_{i1,...,ik}$ is the determinant of k minors of the matrix $A$. However, just like coordinates of points in a projective space, the matrix $A$ is not unique. Since we can multiply on the left any invertible $k \times k$ matrix $\Omega$ without changing the row spaces, because the rows in $\Omega A$ are linear combinations of the rows in $A$.If we choose another basis of $\Lambda$, then the result is multiply by the determinant of $\Omega$. As $\Omega$ invertible, the determinant is never $0$, so the product is always the same point in $\mathbb{P}\BigWedge^{k}V$.\\

Now we try to find these homogeneous polynomials that defines $Im(G(k,n))$. For any k-subspace, we can find k linearly independent vectors that span it. So $\omega \in \mathbb{P}\BigWedge^{k}V$ is in the image if and only if there exits $v_{1},...,v_{k}$(linearly independent) such that $\omega = v_{1}\wedge ... \wedge v_{k}$.

\begin{lm}
    $\omega \in \mathbb{P}\BigWedge^{k}V$ such that $\omega = v_{1}\wedge ... \wedge v_{k}$ if and only if the map:

    \[ V \xrightarrow{\wedge \omega} \mathbb{P}\BigWedge^{k+1}V\]

    has kernel of dimension at least k.
\end{lm}
    
\begin{proof}
    $(\Rightarrow)$ Trivial as ${v_{i}}$ is basis of the kernel.
    $(\Leftarrow)$ If the kernel has dimension at least k, then there exits $v_{1},...,v_{k}$(linearly independent), such that $v_{i}\wedge \omega = 0, 1\leq i \leq k$. \\
    Lets start from $v_{1}$, if $v_{i}\wedge \omega = 0$, then there exists $\omega' \in \mathbb{P}\BigWedge^{k-1}V$ such that $\omega = v_{0}\wedge \omega'$. Repeat this process, we can see that $\omega = v_{1}\wedge ... \wedge v_{k}$.
    
\end{proof}

Then the image can be expressed as:
 \[G = \{\omega \in \mathbb{P}\BigWedge^{k}V \ | \ \text{rank}(V \xrightarrow{\wedge \omega} \mathbb{P}\BigWedge^{k+1}V) \leq n-k  \} \]

This can be interpreted as the zero locus of degree $(n-k+1)$ homogeneous polynomials that are determinant of $(n-k+1)$ minors of the map $\wedge \omega: V \rightarrow \mathbb{P}\BigWedge^{k+1}V$ written as a matrix. So we can conclude that $G(k,n)$ is a variety.



\section{From group of cycles to chow ring}










\newpage
A projective space can be covered by Zariski open subsets isomorphic to affine space. So we can define a local coordinates on projective space. Similarly, we can cover a Grassmannian $G = G(k,n)$ by Zariski open subsets as well. Too see this, fix an $(n-k)$-dimensional subspace $\Gamma \subset V$, and let $U_{\Gamma}$ be the subset of $k$-subspaces that do not meet $\Gamma$:
\[U_{\Gamma} = \{ L\ \in\ G \ | \ L\ \cap\ \Gamma = \emptyset \}
\]


Lets show it is Zariski open: let ${w_1,w_2,\dots ,w_{n-k} }$ be a basis of $\Gamma$ and let
 $\gamma = w_1 \wedge \dots \wedge w_{n-k}$then we have :

\[
U_{\Gamma} = \{ [\omega]   \in  G  \subset   \mathbb{P}(\BigWedge^k V) \ | \ \omega \wedge  \gamma \neq 0 \}
\]



%%%%%%%%%%%%%%%%%%%%%%%%%%%%%%%%%%%%%%%%%%%%%%%%%%%%%%%%%%%%%%%%%%%%%%%%%%%%%%
\newpage

\section{chern classes}

\bibliographystyle{plain}
\bibliography{ref}


\end{document}
